\section{Disclaimer and invitation}
\label{sec:disclaimer}

In fairness, let us be clear: the crypto coma \(\ccoma\) is a \emph{finite} natural number,
however unimaginably large. By calling it a ``cardinality'' and dressing it in the robes of
set theory, we play on a double meaning --- both ``cardinal number'' and ``power/might''.
Therein lies the friendly joke: to clothe the naive dream of the largest number in the
mantle of serious science, while leaving the mathematics itself (tetration, the logarithmic
scale) honest.

And yet ``finite'' does not mean ``comprehensible''. The boys stopped at \(z\) only because
the alphabet ran out, not the numbers; there is no largest number --- every summit turns out
to be a step. Grown-ups ``know'' these numbers --- but they know the \emph{name}, not the
number: no one has ever seen it. Laugh at the two boys once more. It is not quite so funny
now, is it?

\paragraph{An invitation to collaborate.} Ideas, remarks, and your own ascending series are
welcome at \href{mailto:k_mackl@mail.ru?subject=crypto\%20coma}{k\_mackl@mail.ru} and
\href{mailto:ink-shtil@mail.ru?subject=crypto\%20coma}{ink-shtil@mail.ru}, with the subject
line \emph{crypto coma}.
